% Generated draft: requires mathematical and source review.
\documentclass[11pt]{article}
\usepackage[margin=25mm]{geometry}
\usepackage{fontspec}
\setmainfont{Latin Modern Roman}
\usepackage{amsmath,amssymb,mathtools,unicode-math}
\setmathfont{Latin Modern Math}
\usepackage{enumitem,needspace,xurl,hyperref}
\setlength{\parindent}{0pt}
\setlength{\parskip}{5pt}
\setlength{\emergencystretch}{3em}
\newcommand{\N}{\mathbb{N}}
\newcommand{\Z}{\mathbb{Z}}
\newcommand{\Q}{\mathbb{Q}}
\newcommand{\R}{\mathbb{R}}
\newcommand{\C}{\mathbb{C}}
\newcommand{\F}{\mathbb{F}}
\newcommand{\A}{\mathbb{A}}
\newcommand{\PP}{\mathbb{P}}
\newcommand{\E}{\mathbb{E}}
\newcommand{\eps}{\varepsilon}
\newcommand{\dd}{\,\mathrm{d}}
\begin{document}
\section*{UnsolvedMath ID: 30004550}
\textbf{Source problem code:} catalogue entry
\subsection{Definitions and mathematical statement}
Let \(\mathcal B=\N^{\N}\) be Baire space with the product topology, where \(\N\) has the discrete topology. A partial multivalued problem \(f:\subseteq\mathcal B\rightrightarrows\mathcal B\) consists of a domain \(\operatorname{dom}(f)\subseteq\mathcal B\) and a nonempty set \(f(p)\subseteq\mathcal B\) for each \(p\in\operatorname{dom}(f)\). A continuous selector for \(f\) is a continuous map \(F:\operatorname{dom}(f)\to\mathcal B\), with the domain carrying its subspace topology, such that \(F(p)\in f(p)\) for every \(p\in\operatorname{dom}(f)\). Call \(f\) discontinuous if it has no continuous selector.

For two such problems, write \(f\leq_{\mathrm W}^{\mathrm c}g\) if there are partial continuous maps
\[
K:\subseteq\mathcal B\to\mathcal B,
\qquad
H:\subseteq\mathcal B\times\mathcal B\to\mathcal B
\]
satisfying the following requirements. For every \(p\in\operatorname{dom}(f)\), the value \(K(p)\) is defined and belongs to \(\operatorname{dom}(g)\); for every such \(p\) and every \(q\in g(K(p))\), the value \(H(p,q)\) is defined and belongs to \(f(p)\). Thus the output processor may use both the original input and any valid output of \(g\).

Interpreting ``simplest'' as a least continuous Weihrauch degree among all discontinuous problems, the question is whether
\[
\exists f:\subseteq\mathcal B\rightrightarrows\mathcal B\quad
\left[
 f\text{ has no continuous selector}
 \ \land\ 
 \forall g:\subseteq\mathcal B\rightrightarrows\mathcal B,
 \bigl(g\text{ has no continuous selector}
 \Longrightarrow f\leq_{\mathrm W}^{\mathrm c}g\bigr)
\right].
\]
The quantification imposes no effectiveness, definability, or descriptive complexity restriction on the domains or values. The least-degree interpretation is supported by the supplied literature assessment; the word ``simplest'' alone does not distinguish it from the weaker notion of a minimal degree.

The foundational theory for the unrestricted question is not specified. The supplied dated assessment says that Brattka proves that the problem denoted \(\mathrm{DIS}\) is least for effectively discontinuous problems, and that under \(\mathrm{ZF}+\mathrm{DC}+\mathrm{AD}\) it is least among all discontinuous problems. It also reports ``non-effectively-discontinuous counterexamples'' in \(\mathrm{ZFC}\) and describes the unrestricted question as not settled unconditionally. These are supplied literature claims, not independently verified results of this draft. The record does not define \(\mathrm{DIS}\), effective discontinuity, or the precise counterexample property; consequently these reported qualifications are retained in prose rather than expanded into additional formal claims.
\subsection{Short English statement}
Does there exist a discontinuous multivalued problem on Baire space that continuously Weihrauch reduces to every discontinuous multivalued problem on Baire space? Here discontinuity means the absence of a continuous selector, and ``simplest'' is interpreted as least in the reducibility order. The supplied partially solved assessment distinguishes effective discontinuity and results under \(\mathrm{ZF}+\mathrm{DC}+\mathrm{AD}\); its claims require independent review.
\subsection{Sources}
\begin{itemize}
\item Supplied local UnsolvedMath registry record, permanent ID \(30004550\), source problem code \texttt{catalogue entry}, in the collection \emph{Oberwolfach Reports Open Problems}. The record assigns proposed year \(2020\) and supplies no individual proposer. Its labels \texttt{partially\_solved} and \texttt{corrected\_verified} are registry annotations, not independent mathematical verification by this draft.
\item Supplied report citation: \emph{Mathematical Logic: Proof Theory, Constructive Mathematics (2021)}, Oberwolfach Reports, DOI \texttt{10.4171/owr/2020/34}, pp.~1693--1757. \url{https://doi.org/10.4171/owr/2020/34}. The background identifies Oberwolfach Reports as documentation of mathematical workshops and their open-problem discussions. The supplied literature list also labels this reference \emph{Oberwolfach Report 34/2020}.
\item Supplied literature source: Brattka, \emph{The Discontinuity Problem}. \url{https://doi.org/10.1017/jsl.2021.106}.
\item Literature assessment checked on \texttt{2026-08-22}: supplied dated curation derived from the OWR record, not independently verified here. It reports leastness of \(\mathrm{DIS}\) for effectively discontinuous problems, leastness among all discontinuous problems under \(\mathrm{ZF}+\mathrm{DC}+\mathrm{AD}\), and non-effectively-discontinuous counterexamples in \(\mathrm{ZFC}\); it labels the problem partially solved and the unrestricted question not settled unconditionally. The record explicitly requests independent verification of time-sensitive status and cited claims.
\end{itemize}
\end{document}
